\documentclass[11pt]{article}
\usepackage[T1]{fontenc}
\usepackage[utf8]{inputenc}
\usepackage[margin=1in]{geometry}
\usepackage{amsmath,amssymb,amsthm,mathtools}
\usepackage{tikz-cd}
\usepackage{enumitem}
\usepackage[colorlinks=true,linkcolor=blue,citecolor=blue,urlcolor=blue]{hyperref}
\usepackage{microtype}

\newtheorem{theorem}{Theorem}[section]
\newtheorem{proposition}[theorem]{Proposition}
\newtheorem{lemma}[theorem]{Lemma}
\newtheorem{corollary}[theorem]{Corollary}
\theoremstyle{definition}
\newtheorem{definition}[theorem]{Definition}
\newtheorem{example}[theorem]{Example}
\newtheorem{remark}[theorem]{Remark}
\newcommand{\R}{\mathbb R}
\newcommand{\im}{\operatorname{im}}
\newcommand{\rk}{\operatorname{rank}}
\newcommand{\Mat}{\operatorname{Mat}}
\newcommand{\Smooth}{\mathbf{Smooth}}
\newcommand{\Vect}{\mathbf{Vect}_{\R}}
\newcommand{\Lens}{\mathbf{Lens}}
\newcommand{\Para}{\mathbf{Para}}
\newcommand{\Fit}{\operatorname{Fit}}
\newcommand{\Fact}{\mathbf{Fact}}
\newcommand{\id}{\operatorname{id}}
\newcommand{\rev}{\mathcal R}

\title{Finite Evidence and Differentiable Realizations\\
\large A categorical interface between automata learning and gradient learning}
\author{}
\date{Research note \quad 4 November 2026}

\begin{document}
\maketitle

\begin{abstract}
We give a self-contained interface between exact categorical automata learning and
reverse-differentiable learning, using real weighted automata. Finite observations
form a contravariant diagram of Euclidean spaces, while their image hypotheses
form a different, mixed-variance diagram of vector spaces. Polynomial evaluation
of parameterized automata is compatible with observation restriction. Applying
reverse differentiation proves an exact lens-composition law: restriction of
observations is paired with zero-extension of cotangents. An exact-fit,
full-observed-rank realization is isomorphic to the automaton inferred from the
same finite observations. We prove these claims, calculate a two-dimensional
example, and show why rank changes and optimizer updates prevent a stronger
identification. The result is an interface, not a generalization of all of
$\Para(\Lens(\mathcal C))$ or a new active-learning algorithm. It is a limited
technical test for a broader theory of phenomenology, including mathematical
and scientific theory formation; it does not define that broader project.
\end{abstract}

\section{Scope and prior work}

CALF uses wrappers and image factorizations to reconstruct automata from finite
observations, with closedness and consistency governing whether transitions
descend to the inferred state space. Its correctness results additionally use
reachability, observability, and the applicable learning protocol. Its real
weighted instance is presented in $\mathbf{Vect}^{\mathrm{op}}$ and interpreted
in vector spaces; the opposite category accounts for the orientation of its
coalgebraic transitions \cite{calf}.

The image-refinement diagram used below is already present in
Colcombet--Petri\c{s}an--Stabile: their Definition 29 forms the relevant
factorization objects, and Lemma 31 establishes the canonical monomorphisms,
epimorphisms, commuting rectangles, and composition laws. Their FunL$^*$
framework treats automata as functors and covers field-weighted automata
\cite{cps}. We give the elementary finite-dimensional proof because it fixes all
types needed for the differentiable interface. This diagram is not a novelty
claim.

Cruttwell et al. compose parameterized maps, lenses, and reverse differentiation
to describe gradient-based learning. The reverse chain rule supplies a functor
$\rev:\mathcal C\to\Lens(\mathcal C)$ in a Cartesian reverse differential
category \cite{gradient}. We specialize this machinery to the polynomial
evaluation maps of weighted automata.

Two further precedents constrain the interpretation. Spivak already interprets
learners as generalized Moore machines using polynomial functors
\cite{spivak}. Rabusseau, Li, and Precup relate weighted automata to linear
second-order recurrent networks, recover realizations spectrally, and study
gradient-based refinement \cite{rabusseau}. Consequently, neither putting
learners into a dynamical category nor combining spectral identification with
gradient optimization is introduced here.

The results below are an explicit synthesis with proofs. A scoped literature
check is insufficient to establish research novelty, including novelty of the
finite-data compatibility theorem. We distinguish established ingredients,
their proved compatibility in this setting, and questions not resolved by that
compatibility.

\section{Finite observations and their types}

Fix a finite alphabet $A$, its free monoid $A^*$, the empty word $\varepsilon$,
and a target behavior
\[
L:A^*\longrightarrow\R.
\]
The target need not initially be assumed to have finite realization dimension.
All observations in this note are exact real numbers. Approximate oracle
answers require additional theory.

\begin{definition}[Evidence contexts]
An evidence context is a pair $c=(S,E)$ of finite subsets of $A^*$ containing
$\varepsilon$. Let $\mathcal C$ be their poset under componentwise inclusion.
The names ``prefix samples'' and ``suffix tests'' describe their roles; closure
of $S$ under prefixes and of $E$ under suffixes is not assumed for the elementary
matrix results below. Such closure is imposed when an algorithm requires it.
\end{definition}

For a finite set $S$, $\R^S$ denotes its free vector space with basis
$(e_s)_{s\in S}$. Define the target observation and shifted observation maps
\begin{align}
T_c &: \R^S\longrightarrow\R^E,
& (T_c)_{e,s}&=L(se),\label{eq:target}\\
K_{a,c} &: \R^S\longrightarrow\R^E,
& (K_{a,c})_{e,s}&=L(sae)\qquad(a\in A).
\end{align}
Thus columns are prefix samples and rows are suffix tests. These are data,
not parameter-dependent predictions.

Let $\Mat(E,S)=\R^{E\times S}$, with its coordinate Euclidean structure, and set
\[
B_c=\Mat(E,S)\times\prod_{a\in A}\Mat(E,S),
\qquad y_c=(T_c,(K_{a,c})_{a\in A})\in B_c.
\]
For $c=(S,E)\leq d=(S',E')$, write
\[
i_{SS'}:\R^S\hookrightarrow\R^{S'},\qquad
p_{EE'}:\R^{E'}\twoheadrightarrow\R^E
\]
for basis inclusion and coordinate restriction. Define
\begin{equation}\label{eq:restriction}
r_{cd}:B_d\longrightarrow B_c,
\qquad
(X,(X_a)_a)\longmapsto
\bigl(p_{EE'}Xi_{SS'},(p_{EE'}X_ai_{SS'})_a\bigr).
\end{equation}

\begin{proposition}[Observation presheaf]\label{prop:presheaf}
The assignments $c\mapsto B_c$ and $(c\leq d)\mapsto r_{cd}$ define a functor
$B:\mathcal C^{\mathrm{op}}\to\Smooth$. The target data satisfy
$r_{cd}(y_d)=y_c$.
\end{proposition}
\begin{proof}
For $c\leq d\leq f$, the inclusions and projections compose, giving
$r_{cf}=r_{cd}\circ r_{df}$. Identity inclusions and projections give
$r_{cc}=\id$. All these maps are linear, hence smooth. Every retained target
entry is the value of the same concatenated word in both contexts.
\end{proof}

\begin{remark}[Repeated words]
Coordinates of $B_c$ are indexed by $(e,s)$ and $(a,e,s)$, not by distinct words.
Different coordinates can therefore encode the same word. Equation
\eqref{eq:restriction} selects coordinates, so repeated words do not spoil
functoriality. An unweighted Euclidean loss on $B_c$ counts these repetitions
with multiplicity. Alternatively, one may use the finite set
$D_c=SE\cup SAE$ of distinct queried words, its coordinate restriction maps,
and a linear map that duplicates word values into blocks. The convention must
be fixed before choosing a loss.
\end{remark}

\section{Image hypotheses and mixed variance}

Write
\[
I(S,E)=\im(T_{S,E})\subseteq\R^E,
\qquad
\R^S\xrightarrow{\overline T_{S,E}}I(S,E)
\xrightarrow{m_{S,E}}\R^E
\]
for the canonical image factorization. Here $\overline T$ is surjective and
$m$ is inclusion. This is the finite linear instance of a wrapper
factorization. If the target has a realization with state space $Q$, the
wrapper is the sampled reachability map $\R^S\to Q$ followed by restricted
observation $Q\to\R^E$. The image remains defined even when that realization
is unknown.

\begin{proposition}[Mixed image diagram]\label{prop:mixed}
Prefix enlargement at fixed $E$ induces an inclusion
\[
j_{SS'}^E:I(S,E)\hookrightarrow I(S',E).
\]
Suffix enlargement at fixed $S$ induces a surjection in the reverse direction
\[
q_{EE'}^S:I(S,E')\twoheadrightarrow I(S,E).
\]
For $S\subseteq S'$ and $E\subseteq E'$, the square
\[
\begin{tikzcd}
I(S,E')\arrow[r,hook,"j_{SS'}^{E'}"]\arrow[d,two heads,"q_{EE'}^S"']&
I(S',E')\arrow[d,two heads,"q_{EE'}^{S'}"]\\
I(S,E)\arrow[r,hook,"j_{SS'}^E"']&I(S',E)
\end{tikzcd}
\]
commutes. Both kinds of arrows preserve identities and compose in the evident
way. Thus $I$ is a strict functor in $S$ covariantly and $E$ contravariantly.
\end{proposition}
\begin{proof}
At fixed $E$, enlarging $S$ adds columns to $T$, so the old image is a subspace
of the new image. At fixed $S$,
$p_{EE'}T_{S,E'}=T_{S,E}$, so restriction maps the former image onto the latter.
Both composites in the square take a vector to its restriction to $E$ in
$\R^E$. Inclusions and coordinate restrictions compose strictly.
\end{proof}

This is the vector-space specialization of \cite[Definition 29 and Lemma
31]{cps}. It can also be read as a double diagram: horizontal maps are
injections, vertical maps are surjections, and two-dimensional cells are
commuting squares. Horizontal and vertical composition are ordinary pasting;
the interchange law follows from composition of linear maps. No failure of
functoriality is being asserted. Nor is a commuting square automatically a
pullback square.

The two diagrams $B$ and $I$ have different variance. Increasing both samples
and tests restricts the available observation \emph{coordinates} from $B_d$ to
$B_c$, but on image \emph{state spaces} the two operations induce arrows in
opposite directions. This distinction is essential when discussing refinement.

\section{When the observed dynamics descend}

Fix a context and abbreviate $T=T_c$, $K_a=K_{a,c}$, and $I=\im(T)$.

\begin{proposition}[Linear closedness and consistency]\label{prop:descent}
There exists a unique linear transition $U_a:I\to I$ satisfying
\begin{equation}\label{eq:descent}
mU_a\overline T=K_a
\end{equation}
if and only if
\begin{equation}\label{eq:cc}
\im(K_a)\subseteq\im(T),\qquad \ker(T)\subseteq\ker(K_a).
\end{equation}
\end{proposition}
\begin{proof}
Necessity follows by factoring $K_a$ as in \eqref{eq:descent}. Conversely,
define $U_a(Tv)=K_av$, interpreting the right-hand side as an element of $I$.
The image condition makes the codomain correct. If $Tv=Tw$, then
$v-w\in\ker(T)\subseteq\ker(K_a)$, so the definition is independent of the
representative. It is linear. Since $\overline T$ is surjective, the defining
equation determines $U_a$ uniquely.
\end{proof}

When these conditions hold for every $a$, define initial state and output by
\[
\alpha_I=T e_\varepsilon\in I,
\qquad o_I:I\to\R,\quad o_I(x)=x_\varepsilon.
\]
This produces an inferred weighted automaton. The presence of $\varepsilon$
in both indexing sets supplies these two maps without further choices.
The conditions \eqref{eq:cc} are the image and kernel forms of the relevant
CALF closedness and consistency requirements \cite{calf}.

\begin{remark}[Block dynamics need not reproduce arbitrary word samples]
For arbitrary finite $S,E$, Proposition~\ref{prop:descent} proves an algebraic
factorization of the shifted blocks. It does not by itself prove that the
resulting initial state and output reproduce all sampled words. For example,
take $A=\{a\}$, $S=\{\varepsilon,aa\}$, $E=\{\varepsilon\}$ and target
values $L(\varepsilon)=1$, $L(a)=3$, $L(aa)=2$, $L(aaa)=6$. Then
$T=[1\;2]$, $K=[3\;6]$ satisfy both image and kernel conditions, but the
descended automaton has transition $3$, initial value $1$, and identity
output. It predicts $9$, not $2$, on $aa$. The missing intermediate prefix
$a$ matters. The next proposition supplies sufficient closure conditions.
Alternatively, the exact-fit hypothesis of
Theorem~\ref{thm:compatibility} supplies reproduction through the candidate
realization itself.
\end{remark}

\begin{proposition}[Reproduction under word closure]\label{prop:reproduction}
Suppose $S$ is prefix-closed, $E$ is suffix-closed, both contain
$\varepsilon$, and the blocks are coherent word observations as in
\eqref{eq:target}. If the conditions \eqref{eq:cc} hold for every $a$, the
inferred automaton reproduces $L$ on $SE\cup SAE$.
\end{proposition}
\begin{proof}
Use $U_{a_1\cdots a_k}=U_{a_k}\cdots U_{a_1}$. Prefix closure and induction
give $U_s\alpha_I=T e_s$ for each $s\in S$: if $s=ta$, then
$U_aT e_t=K_a e_t=T e_{ta}$, as vectors indexed by $E$.

Suffix closure gives $o_IU_e x=x_e$ for all $x\in I$ and $e\in E$.
For $e=\varepsilon$ this is the definition. Write $e=av$, with $v\in E$,
and $x=Tz$. By induction,
\[
o_IU_e x=o_IU_vU_a x=(U_a x)_v
=(K_a z)_v=\sum_{s\in S}z_sL(sav)=(Tz)_{av}=x_e.
\]
Consequently $o_IU_eU_s\alpha_I=(T e_s)_e=L(se)$ and
$o_IU_eU_aU_s\alpha_I=(K_a e_s)_e=L(sae)$.
\end{proof}

\begin{proposition}[Compatibility of valid adjacent hypotheses]\label{prop:intertwine-grid}
If both endpoints of a horizontal or vertical arrow in
Proposition~\ref{prop:mixed} admit descended transitions, that arrow is a
homomorphism of their inferred automata.
\end{proposition}
\begin{proof}
For a vertical arrow $q:I(S,E')\to I(S,E)$, coordinate restriction gives
$q\overline T'=\overline T$ and maps the corestriction of $K'_a$ to that of
$K_a$. Hence
\[
qU'_a\overline T'=q\overline K'_a
=\overline K_a=U_a\overline T=U_aq\overline T'.
\]
Cancel the surjection $\overline T'$. Initial states are preserved by the same
restriction identity, and the $\varepsilon$ coordinate is retained. For a
horizontal inclusion $j$, use
$j\overline T=\overline T'i_{SS'}$ and the identical relation for shifted
blocks. Then $U'_aj\overline T=jU_a\overline T$; cancel $\overline T$.
Initial state and output are again preserved.
\end{proof}

This assertion applies only where both hypotheses exist. A refinement needed
to repair a defect passes through a cell at which a transition may not descend.
It does not make successive, disagreeing conjectures globally equivalent.

\section{Differentiable realization charts}

For $n\geq1$, let
\[
\Theta_n=\R^n\times\prod_{a\in A}\R^{n\times n}\times\R^n.
\]
Write $\theta=(\alpha,(M_a)_a,\beta)$ and use the convention
\[
M_\varepsilon=I_n,\qquad
M_{a_1\cdots a_k}=M_{a_k}\cdots M_{a_1},\qquad
L_\theta(w)=\beta^{\mathsf T}M_w\alpha.
\]
Thus $M_{uv}=M_vM_u$, so $L_\theta(sae)=\beta^{\mathsf T}M_eM_aM_s\alpha$.
Replacing $L$ by $L_\theta$ in \eqref{eq:target} defines predicted blocks
$\widehat T_{\theta,c},\widehat K_{a,\theta,c}$ and a polynomial evaluator
\[
\Phi_{n,c}:\Theta_n\longrightarrow B_c,
\qquad
\theta\longmapsto
(\widehat T_{\theta,c},(\widehat K_{a,\theta,c})_a).
\]
Words index the family of polynomial computations. We do not regard the
unbounded discrete set $A^*$ as a Euclidean object of $\Smooth$.

\subsection{The relevant lens operations}

An object of $\Lens(\Smooth)$ is a pair $(X,X')$. A lens
$(f,f^*):(X,X')\to(Y,Y')$ consists of
$f:X\to Y$ and $f^*:X\times Y'\to X'$. Composition is
\[
(g,g^*)\circ(f,f^*)
=\bigl(g\circ f,(x,z')\mapsto f^*(x,g^*(f(x),z'))\bigr).
\]
For a smooth map $f:X\to Y$ between Euclidean spaces,
\[
\rev(f)=(f,R[f]),\qquad R[f](x,v)=Jf(x)^{\mathsf T}v.
\]
The identity and chain rules make $\rev$ a functor. The Euclidean inner products
identify cotangent vectors with coordinate vectors; this does not mean
cotangents have the same semantic role as forward values.

A morphism in $\Para(\Lens(\Smooth))$ adds a parameter pair $(P,P')$ to such
a lens. Therefore $\Phi_{n,c}$ defines a parametric lens from $(1,1)$ to
$(B_c,B_c)$ with parameter pair $(\Theta_n,\Theta_n)$, where $1=\R^0$:
its backward computation is $(\theta,v)\mapsto J\Phi_{n,c}(\theta)^{\mathsf T}v$.

\begin{theorem}[Observation restriction and reverse differentiation]\label{thm:bridge}
For every $c\leq d$ and every fixed $n$,
\begin{align}
\Phi_{n,c}&=r_{cd}\circ\Phi_{n,d},\label{eq:forwardbridge}\\
\rev(\Phi_{n,c})&=\rev(r_{cd})\circ\rev(\Phi_{n,d}),\label{eq:lensbridge}\\
R[\Phi_{n,c}](\theta,v)
&=R[\Phi_{n,d}](\theta,r_{cd}^{\mathsf T}v).\label{eq:backbridge}
\end{align}
Here $r_{cd}^{\mathsf T}:B_c\to B_d$ extends indexed coordinate vectors by
zero. These identities compose consistently along all chains of refinements.
\end{theorem}
\begin{proof}
Each retained predicted entry is given by the same word polynomial in the two
contexts, proving \eqref{eq:forwardbridge}. The Jacobian of $r_{cd}$ is the
constant coordinate projection $r_{cd}$. Transposing the chain rule gives
\eqref{eq:backbridge}, which is exactly the backward part of
\eqref{eq:lensbridge}. For $c\leq d\leq f$, Proposition~\ref{prop:presheaf}
gives $r_{cf}^{\mathsf T}=r_{df}^{\mathsf T}r_{cd}^{\mathsf T}$, so the two
routes agree. Identities follow in the same way.
\end{proof}

\begin{corollary}[Masked loss compatibility]\label{cor:masked}
For $\ell_c(\theta)=\tfrac12\|\Phi_{n,c}(\theta)-y_c\|^2$,
\[
\nabla\ell_c(\theta)=J\Phi_{n,d}(\theta)^{\mathsf T}
r_{cd}^{\mathsf T}r_{cd}\bigl(\Phi_{n,d}(\theta)-y_d\bigr).
\]
The full $d$-loss gradient equals this vector plus the contribution of the
new coordinates.
\end{corollary}
\begin{proof}
Substitute \eqref{eq:forwardbridge} and $y_c=r_{cd}y_d$ into the derivative of
the squared Euclidean norm. The matrix $r_{cd}^{\mathsf T}r_{cd}$ is the
orthogonal coordinate mask. Its complement supplies the other summand.
\end{proof}

The local primitive behind this construction is
$F(M,x)=Mx$, whose reverse derivative is
\[
R[F](M,x,v)=(vx^{\mathsf T},M^{\mathsf T}v).
\]
Lens composition handles sequential transitions. When a symbol occurs several
times, a diagonal reparameterization shares its matrix; the reverse map of
that diagonal adds its separate gradient contributions. Thus the construction
uses composition, not merely an independently differentiated final formula.

\section{Exact fit and the factorization category}

\begin{definition}[Factorizations of one observed matrix]
For $T:\R^S\to\R^E$, the category $\Fact(T)$ has objects
\[
\R^S\xrightarrow{R}V\xrightarrow{O}\R^E,
\qquad OR=T,
\]
where $V$ is finite-dimensional. A morphism to $(V',R',O')$ is a linear map
$h:V\to V'$ such that $hR=R'$ and $O'h=O$. Composition and identities are
ordinary linear composition and identities.
\end{definition}

\begin{proposition}[Universal properties of the observed image]\label{prop:universal}
The image factorization is terminal in the full subcategory where $R$ is
surjective and initial in the full subcategory where $O$ is injective.
\end{proposition}
\begin{proof}
If $R$ is surjective, $\im(O)=\im(OR)=\im(T)$. Corestricting $O$ therefore
gives the unique map $h:V\to\im(T)$ with $mh=O$ and
$hR=\overline T$. If $O$ is injective, define
$h:\im(T)\to V$ by $h(Tv)=Rv$. If $Tv=Tw$, injectivity of $O$ implies
$Rv=Rw$. This is a linear map satisfying $h\overline T=R$ and $Oh=m$.
Surjectivity of $\overline T$ proves uniqueness.
\end{proof}

Define the exact-fit set
\[
\Fit_n(c)=\{\theta\in\Theta_n:\Phi_{n,c}(\theta)=y_c\}.
\]
Theorem~\ref{thm:bridge} makes $\Fit_n$ a presheaf on $\mathcal C$ with
restriction given by inclusion $\Fit_n(d)\subseteq\Fit_n(c)$.
Its category of elements has objects $(c,\theta)$ and an arrow
$(c,\theta)\to(d,\theta')$ when $c\leq d$ and $\theta'=\theta$ survives the
additional equations. This records refinement of exact evidence. Gradient
steps need not lie in this category, since intermediate parameters generally
fail exact-fit equations.

For any candidate define
\begin{align*}
R_{\theta,S}:\R^S&\to\R^n,
&R_{\theta,S}e_s&=M_s\alpha,\\
O_{\theta,E}:\R^n&\to\R^E,
&(O_{\theta,E}x)_e&=\beta^{\mathsf T}M_ex.
\end{align*}
Then always
\begin{equation}\label{eq:candidatefactor}
\widehat T_{\theta,c}=O_{\theta,E}R_{\theta,S},
\qquad
\widehat K_{a,\theta,c}=O_{\theta,E}M_aR_{\theta,S}.
\end{equation}
These factor the \emph{target} blocks only when exact agreement holds.

\begin{theorem}[Finite-data realization compatibility]\label{thm:compatibility}
Let $T_c,(K_{a,c})_a$ be target observations and suppose
$\theta\in\Fit_n(c)$ with $\rk(T_c)=n$. Then the target blocks satisfy
\eqref{eq:cc}; the inferred automaton exists; and the corestriction
\[
\overline O:\R^n\longrightarrow I_c=\im(T_c)
\]
of $O_{\theta,E}$ is an isomorphism of the candidate automaton with that
inferred automaton. In particular,
\[
\overline OM_a=U_a\overline O,
\qquad \overline O\alpha=\alpha_I,
\qquad o_I\overline O=\beta^{\mathsf T}.
\]
\end{theorem}
\begin{proof}
Abbreviate $R=R_{\theta,S}$ and $O=O_{\theta,E}$. Exact fit gives
$T_c=OR$ and $K_{a,c}=OM_aR$. Since $\rk(T_c)=n$ and the intermediate space
has dimension $n$, $R$ is surjective and $O$ injective. Thus
$\im(O)=\im(T_c)$, making $\overline O$ an isomorphism, and
$\ker(T_c)=\ker(R)\subseteq\ker(K_{a,c})$. Also
$\im(K_{a,c})\subseteq\im(O)=\im(T_c)$. Proposition~\ref{prop:descent}
therefore supplies $U_a$. Corestricting the defining equations gives
\[
\overline OM_aR=\overline K_{a,c}
=U_a\overline T_c=U_a\overline OR.
\]
Cancel the surjection $R$ to obtain the transition identity. Since
$R e_\varepsilon=\alpha$, the initial-state identity follows. Since
$M_\varepsilon=I_n$, the $\varepsilon$ coordinate of $Ox$ is
$\beta^{\mathsf T}x$, proving the output identity.
\end{proof}

\begin{corollary}[A conditional dimension certificate]
Every candidate matching $T_c$ has dimension at least $\rk(T_c)$. Any two
exact-fit candidates of dimension exactly $\rk(T_c)$ are isomorphic as
automata if they also match all the shifted blocks.
\end{corollary}
\begin{proof}
The rank bound follows from $T_c=OR$. The second assertion follows by
composing the two isomorphisms of Theorem~\ref{thm:compatibility} through $I_c$.
\end{proof}

\paragraph{Boundary of this theorem.}
It identifies the candidate with the automaton inferred from the finite table.
It does not prove that either agrees with the unknown target on every word.
For that conclusion one needs an appropriate complete sampling condition,
known bounds with a complete test, or a sound equivalence oracle. A sufficient
concrete condition is that the target itself has dimension $n$ and the
sampled reachability and observation maps have full rank; applying the same
argument to the target identifies it with $I_c$. Finite rank of one block
alone does not establish that condition.

\section{A complete two-dimensional example}

Let $A=\{a\}$ and
\[
L(a^k)=1+2^k\quad(k\geq0).
\]
The first values are $2,3,5,9,17$. A realization is
\[
\alpha=\beta=\begin{pmatrix}1\\1\end{pmatrix},
\qquad M=\begin{pmatrix}1&0\\0&2\end{pmatrix}.
\]

\subsection{Refinement and its two distinct defects}
For $S=E=\{\varepsilon\}$,
\[
T=[2],\qquad K=[3].
\]
The inferred one-state automaton has transition $3/2$, initial value $2$, and
output the identity. It predicts $2(3/2)^k$, hence predicts $9/2$ rather than
$5$ at $a^2$. An exact test of $a^2$ therefore refutes it. An equivalence
oracle could return that word, but merely constructing the table does not
provide such an oracle.

If prefixes are enlarged first to $S'=\{\varepsilon,a\}$, then
\[
T=[2\;3],\qquad K=[3\;5].
\]
The vector $(3,-2)^{\mathsf T}$ lies in $\ker(T)$ but is sent to $-1$ by $K$.
Consistency fails. If suffixes are enlarged first to
$E'=\{\varepsilon,a\}$, then
\[
T=\begin{pmatrix}2\\3\end{pmatrix},
\qquad K=\begin{pmatrix}3\\5\end{pmatrix}.
\]
The second vector is not in the span of the first, so closedness fails.

With both enlarged,
\[
T=\begin{pmatrix}2&3\\3&5\end{pmatrix},
\quad
K=\begin{pmatrix}3&5\\5&9\end{pmatrix},
\quad
T^{-1}=\begin{pmatrix}5&-3\\-3&2\end{pmatrix}.
\]
Now $\det(T)=1$ and $I=\R^2$. The canonical observed-state automaton is
\[
U=KT^{-1}=\begin{pmatrix}0&1\\-2&3\end{pmatrix},
\qquad \alpha_I=\begin{pmatrix}2\\3\end{pmatrix},
\qquad o_I=\begin{pmatrix}1&0\end{pmatrix}.
\]
It implements $L(a^{k+2})=3L(a^{k+1})-2L(a^k)$.

The image rectangle is explicitly
\[
\begin{tikzcd}
\operatorname{span}\{(2,3)^{\mathsf T}\}
\arrow[r,hook]\arrow[d,two heads,"\mathrm{pr}_1"']&
\R^2\arrow[d,two heads,"\mathrm{pr}_1"]\\
\R\arrow[r,equal]&\R.
\end{tikzcd}
\]
Thus the categorical rectangle commutes even though its two intermediate
cells do not both support valid inferred transitions.

\subsection{Exact compatibility with the differentiable realization}
For the diagonal realization,
\[
R=\begin{pmatrix}1&1\\1&2\end{pmatrix},
\qquad
O=\begin{pmatrix}1&1\\1&2\end{pmatrix},
\qquad T=OR,\qquad K=OMR.
\]
Direct multiplication gives
\[
OM=\begin{pmatrix}1&2\\1&4\end{pmatrix}=UO.
\]
The matrix $O$ sends $\alpha$ to $(2,3)^{\mathsf T}$ and its first row is
$\beta^{\mathsf T}$, verifying every part of
Theorem~\ref{thm:compatibility}.

A one-state realization has behavior $cr^k$ with $c=\alpha\beta$.
Matching the first two values forces $c=2$ and $r=3/2$, contradicting the
third value. The rank-two block certifies the obstruction independently of
optimization. The displayed two-state parameters achieve zero loss on the
observed blocks. The fact that a zero-loss point exists does not prove that
gradient descent reaches it.

\section{Why this is an interface rather than an identification}

\subsection{Rank changes obstruct a single smooth minimal-state bundle}

\begin{proposition}[Rank-jump obstruction]\label{prop:rankjump}
There is no vector bundle over an interval containing $0$ whose fiber at $t$
is isomorphic to $\im(H(t))$ for every
\[
H(t)=\begin{pmatrix}1&0\\0&t\end{pmatrix}.
\]
Consequently, exact minimal image extraction for arbitrary smooth matrix
families cannot always be represented by a fixed-rank smooth state bundle.
\end{proposition}
\begin{proof}
At $t=0$ the image has dimension one; at every nonzero $t$ it has dimension
two. A vector bundle has locally constant fiber dimension, giving a
contradiction in every neighborhood of zero.
\end{proof}

This is not a no-go theorem for stratified spaces, variable-dimensional
semantics, padded state spaces, or more general categories. It identifies the
extra structure needed beyond a fixed Euclidean chart. It also explains a
numerical warning: $H(t)^{-1}=\operatorname{diag}(1,1/t)$ for $t\ne0$, so
inverse-based coordinates need not remain bounded near a rank drop.

\subsection{The lens law does not make gradient updates natural}

Theorem~\ref{thm:bridge} compares derivatives with the correct transformed
cotangent seed. It does not identify the full losses before and after adding
evidence. Likewise, a behavior-preserving change of coordinates need not
preserve a Euclidean gradient step. For a scalar realization with prediction
$\alpha\beta$, consider
\[
\ell(\alpha,\beta)=\tfrac12(\alpha\beta)^2,
\qquad g_c(\alpha,\beta)=(c\alpha,\beta/c).
\]
The loss is invariant under $g_c$, but at $(1,1)$ its gradient is $(1,1)$.
For $c=2$, the gradient at $g_2(1,1)=(2,1/2)$ is $(1/2,2)$, whereas
$Dg_2|_{(1,1)}[(1,1)]=(2,1/2)$. Thus pushing a gradient update through this realization
isomorphism gives a different update. A choice of metric or compatible
optimizer geometry is additional data.

\subsection{An observed quotient can preserve dynamics but lose gradient form}

Consider $f_p=p^2$ with target $b$ and squared loss
$\ell(p)=\tfrac12(p^2-b)^2$. A gradient step is
\[
F_{b,\eta}(p)=p-2\eta p(p^2-b).
\]
The readout $q(p)=p^2$ yields an exact quotient:
\[
q\circ F_{b,\eta}=\overline F_{b,\eta}\circ q,
\qquad
\overline F_{b,\eta}(q)=q\bigl(1-2\eta(q-b)\bigr)^2,
\quad q\geq0.
\]
The equality follows by squaring $F_{b,\eta}(p)$. For a Moore machine whose
inputs are the chosen pairs $(b,\eta)$ and visible output is $p^2$, the relation
$p\sim p'$ iff $p^2=(p')^2$ is a behavioral congruence: updates are odd and
therefore preserve sign-related pairs. Its quotient generally has infinitely
many states. The quotient update is not ordinary Euclidean gradient descent
on $\tfrac12(q-b)^2$: expanding it gives
$q-4\eta q(q-b)+4\eta^2q(q-b)^2$. This example concerns observed training
dynamics, not direct identification of an unknown predictor. It illustrates
why observational quotienting and differentiable learning are different
operations.

\section{What has been constructed and what is still required}

The mathematical object assembled here consists of the observation presheaf
$B$, the target section $y$, the mixed image diagram $I$, polynomial candidate
cones $\Phi_{n,-}$, their reverse-derivative lens images, and categories of
factorizations and exact-fit realizations. Its composition laws and the
interface theorem are explicit. On valid image cells it also carries inferred
automata and intertwining maps. This is more than an analogy, but does not
identify all these morphisms with each other.

Fixing $c,n$ recovers ordinary reverse-differentiable fitting of one weighted
automaton chart. Forgetting the charts recovers finite image hypotheses and
their exact refinement data. At a full-observed-rank exact fit,
Theorem~\ref{thm:compatibility} identifies the two descriptions of the
finite-data hypothesis. These are recovered cases, not a claim that arbitrary
CALF instances are reverse differential categories or that all parametric
lenses are automata learners.

A full combined learning theory still needs:
\begin{enumerate}[leftmargin=*,itemsep=3pt]
\item A teacher or experiment interface and an authorized rule for choosing
new samples, tests, or counterexamples. An update of parameters is not itself
an equivalence query.
\item A termination or convergence argument with stated finiteness,
observability, optimization, and noise assumptions. The two original
frameworks do not supply these assumptions for one another.
\item A semantics for changing realization dimension and preserving useful
witnesses when a quotient or rank changes.
\item An appropriate differential or metric structure on realization
isomorphism classes if coordinate-independent optimizer behavior is wanted.
\end{enumerate}
Witness-based automata learning is already developed in the apartness-based
$L^{\#}$ approach \cite{apartness}; learners and teachers themselves have also
been treated coalgebraically for formalized complexity analysis
\cite{coalgebraic}. Those are relevant existing structures, not gaps created
by an absence of categorical work on interaction or evidence.

\section{Implications for a broader theory of phenomenology}

The broader motivation includes how partial characterizations support the
formation and revision of mathematical and scientific theories. This note
offers one technical test of that motivation. It does not define phenomenology
by automata learning, observation functions, or gradient descent.

The finite context $(S,E)$ determines a partial characterization of the target.
The image $I(S,E)$ is structure generated from that characterization through
an explicit universal factorization. The refinement diagram records how that
generated structure changes when samples or tests change. A differentiable
realization supplies coordinates and a mechanism for adjusting a candidate
description. The exact-fit theorem states when such a description and the
generated structure agree up to isomorphism. These are precise statements in
the weighted-automaton setting, not assertions that every scientific or
mathematical theory has this form.

The construction separates evidence, generated structure, and transport.
Restriction transports observation coordinates. Mixed-variance maps transport
image hypotheses. The reverse derivative transports covectors through a
chosen realization. Their relations are proved by the bridge law and the
finite-data compatibility theorem; their differences remain visible in their
types. In particular, an update within an existing description is different
from an enlargement of the distinctions that a description can represent.
The rank-jump example makes this difference mathematically unavoidable in a
fixed-rank smooth setting.

For a broader application one would have to specify what counts as evidence,
which presentations and equivalences are relevant, how new tests or concepts
are generated, and what justifies a claim beyond the finite characterization.
The present results neither supply those choices universally nor show that
their mechanisms are shared by all domains. Their contribution is a worked
instance in which partial characterization, generated structure, and
differentiable adjustment can be compared without identifying them, together
with explicit conditions under which information passes between them.

\begin{thebibliography}{9}

\bibitem{calf}
Gerco van Heerdt, Matteo Sammartino, and Alexandra Silva.
\emph{CALF: Categorical Automata Learning Framework}.
CSL 2017; arXiv:1704.05676, version 3.
\url{https://arxiv.org/abs/1704.05676}.
The weighted-automaton construction is in Appendix A of this version.

\bibitem{gradient}
G.~S.~H. Cruttwell, Bruno Gavranovi\'{c}, Neil Ghani, Paul Wilson, and Fabio Zanasi.
\emph{Categorical Foundations of Gradient-Based Learning}.
arXiv:2103.01931, version 2, 2021.
\url{https://arxiv.org/abs/2103.01931}.

\bibitem{cps}
Thomas Colcombet, Daniela Petri\c{s}an, and Riccardo Stabile.
\emph{Learning Automata and Transducers: A Categorical Approach}.
CSL 2021; arXiv:2010.13675.
\url{https://arxiv.org/abs/2010.13675}.
Definition 29 and Lemma 31 in the extended version:
\url{https://www.irif.fr/~colcombe/Publications/CSL21-colcombet-petrisan-stabile-long.pdf}.

\bibitem{spivak}
David I. Spivak.
\emph{Learners' Languages}.
Electronic Proceedings in Theoretical Computer Science 372 (2022), 14--28.
arXiv:2103.01189; version 3, 2025.
\url{https://arxiv.org/abs/2103.01189}.
\href{https://doi.org/10.4204/EPTCS.372.2}{doi:10.4204/EPTCS.372.2}.

\bibitem{rabusseau}
Guillaume Rabusseau, Tianyu Li, and Doina Precup.
\emph{Connecting Weighted Automata and Recurrent Neural Networks through
Spectral Learning}.
Proceedings of AISTATS, PMLR 89 (2019), 1630--1639.
\url{https://proceedings.mlr.press/v89/rabusseau19a.html}.

\bibitem{apartness}
Frits Vaandrager, Bharat Garhewal, Jurriaan Rot, and Thorsten Wi\ss mann.
\emph{A New Approach for Active Automata Learning Based on Apartness}.
arXiv:2107.05419, version 4, 2022.
\url{https://arxiv.org/abs/2107.05419}.

\bibitem{coalgebraic}
Thorsten Wi\ss mann.
\emph{Formalized Run-Time Analysis of Active Learning --- Coalgebraically in Agda}.
arXiv:2602.16427, version 2, 2026.
\url{https://arxiv.org/abs/2602.16427}.

\end{thebibliography}
\end{document}
